\chapter*{The internship in two pages}
\addcontentsline{toc}{chapter}{The internship in two pages}
\thispagestyle{plain}

\paragraph{The goal.} Predict which antenna a mobile phone will be attached to
next, from the beginning of that person's day and nothing else. An operator wants
this because it decides where capacity is needed, and because a connection
prepared before it is required is a connection that does not drop.

\paragraph{The data.} Fifteen daily exports from a Czech operator covering the
Karlovy Vary area, 12 to 26 March 2014. 369 antennas mounted on 113 masts, and
one line per person per day: an anonymous number, then the antennas that person's
phone used with the time of each. No names, no numbers, no content. A day file
holds between 80{,}217 and 99{,}504 people.

\paragraph{The competitor.} The rule \emph{``the person is still on the same
antenna as a moment ago''}. One line of code, no training. On a day file taken as
it comes it is right 69.2 times out of a hundred, because people spend most of
their time not going anywhere. Everything in this report is a distance to that
rule, never an absolute score.

\paragraph{What I built.} A small general-purpose language model, retrained on
those days until it stops being a writer of English and becomes a predictor of
movement in one Czech region. Every run was trained on one free graphics card.
The whole instrument was fixed in the first week and never changed afterwards:
what changed, eleven times, was the data and the training signal.

\paragraph{What I actually did, week by week.} Table~\ref{tab:recap} is the
internship. Each row is one iteration: which people it was measured on, the one
thing I changed that week, and how far the trained model landed from the one-line
rule.

% T14 -- Every headline run of the internship on one axis, gap by gap.
\begin{table}[t]
\centering\footnotesize
\caption{Every iteration, on one axis}
\label{tab:recap}
\begin{tabular}{@{}c@{\hspace{8pt}}l@{\hspace{8pt}}l@{\hspace{8pt}}l@{\hspace{8pt}}r@{\hspace{8pt}}r@{\hspace{9pt}}r@{}}
\toprule
\textbf{It.} & \textbf{Corpus} & \textbf{Lever of that iteration} & \textbf{Protocol} & \textbf{Base.} & \textbf{Model} & \textbf{Gap} \\
\midrule
1  & \sraw     & pipeline built    & last 20\,\% & 69.2\,\% & 16.9\,\% & $-52.3$ \\
2  & \sraw     & harness fixes     & last 20\,\% & 69.2\,\% & 64.4\,\% & $-4.8$  \\
3  & \srawbig  & ten times more users & last 20\,\% & 81.6\,\% & 76.0\,\% & $-5.6$  \\
4  & \sfilt    & filtering         & last 20\,\% & 79.0\,\% & 80.1\,\% & $\mathbf{+1.1}$ \\
5  & \smob     & a cleaner corpus  & last 20\,\% & 45.5\,\% & 44.6\,\% & $-0.9$  \\
6  & \smob     & hour tokens       & last 20\,\% & 45.5\,\% & 44.2\,\% & $-1.3$  \\
7  & \smob     & group token       & full seq.   & 43.1\,\% & 41.3\,\% & $-1.8$  \\
8  & \smob     & adapter routing   & full seq.   & 43.1\,\% & 46.8\,\% & $\mathbf{+3.7}$ \\
9  & \stwok    & 1{,}900 users     & last 20\,\% & 48.2\,\% & 52.4\,\% & $\mathbf{+4.2}$ \\
9  & \slev     & levelled records  & last 20\,\% & 76.5\,\% & 74.9\,\% & $-1.6$  \\
10 & \seleven  & 21{,}000 users    & last 20\,\% & 50.2\,\% & 54.5\,\% & $\mathbf{+4.3}$ \\
11 & \seleven  & analysis only     & \multicolumn{4}{@{}l@{}}{\textit{no new training: why the ladder climbs, then stops}} \\
\bottomrule
\end{tabular}

\figtext{%
The eleven iterations reduced to the only quantity that is comparable between
them, the gap in points between the model and the persistence baseline of its
own test users. Reading down the \textbf{Gap} column is reading the whole
internship: five iterations below the trivial rule, one accidental win at
iteration~4, and a real one from iteration~8 onward. Reading down the
\textbf{Base.} column shows why no other column may be compared vertically,
since the corpus, and with it the difficulty of the test users, changed four
times. What each iteration actually did is in
Appendix~\ref{app:timeline}.}

\vspace{4pt}
\begin{minipage}{\textwidth}\footnotesize
``Last 20\,\%'' scores the final fifth of each test sequence, ``full seq.''
every position of it, one at a time, which is strictly harder and lowers the
baseline on the very same users from 45.5\,\% to 43.1\,\%
(Chapter~\ref{ch:time}). Iteration~6 is the row whose single number understates
it: its targeted sampling won $+2.6\pt$ inside the two guessed windows and lost
more than that outside them. Iteration~9 appears twice because two independent
things were tried that week, more users and levelled records.
\end{minipage}
\end{table}


\paragraph{What came out of it.} Three findings, each established in the chapter
named.

\begin{enumerate}
  \item \textbf{The model beats the rule by 4.3 more correct answers per hundred},
        54.5\,\% against 50.2\,\% on a thousand people it had never seen, and by
        4.2 on a second, separate group. What produced that gain was neither a
        bigger model nor more data, but sorting people by \emph{the hours at which
        they actually change location} and giving each kind of person its own set
        of trained numbers (Chapter~\ref{ch:groups}).
  \item \textbf{Population size stops paying at about two thousand people.} The
        edge grows from 0.6 to 4.2 correct answers per hundred between 400 and
        1{,}900 training people, and a further ten-fold increase to 21{,}000 buys
        0.1, with zero inside every confidence interval
        (Chapter~\ref{ch:scale}).
  \item \textbf{Difficulty lives in the person, not in the dataset.} 22.4 correct
        answers per hundred separate the calmest kind of user from the most
        restless one, against 0.4 across the entire range of training sizes: a
        factor of more than fifty (Chapters~\ref{ch:user} and~\ref{ch:scale}).
\end{enumerate}

\paragraph{And one finding about measurement rather than modelling.} The three
largest single improvements of the nine weeks were not modelling ideas. Reading
the best saved copy of a run instead of the last one was worth 15.5 correct
answers per hundred. Recognising that where a day is cut in two belongs to the
measurement and not to the model killed an apparent win. And choosing people on
whom the one-line rule does not already score 69\,\% is what made the task
measurable at all. On a problem whose honest yardstick is one line of code, the
measurement deserves to be audited before the model is
(Chapter~\ref{ch:tools}).
